\documentclass[11pt,a4paper]{article}
\usepackage[utf8]{inputenc}
\usepackage[vietnamese]{babel}
\usepackage{amsmath,amssymb}
\usepackage{geometry}
\geometry{a4paper, total={180mm,260mm}, left=15mm, right=15mm, top=18mm, bottom=18mm}

\usepackage{tikz}
\usetikzlibrary{shapes.geometric, arrows.meta, positioning, calc}
\tikzset{
  arr/.style={-{Stealth[scale=0.8]}},
  every node/.append style={align=center}
}
\usepackage{tcolorbox}
\tcbuselibrary{skins,breakable}
\usepackage{fontawesome}
\usepackage{fancyhdr}
\usepackage{enumitem}

\pagestyle{fancy}
\fancyhf{}
\renewcommand{\headrulewidth}{0.5pt}
\renewcommand{\footrulewidth}{0.5pt}
\fancyhead[L]{\small\bfseries\color{violet!85!black}\faBook\ SỔ TAY KIẾN THỨC TOÁN 2}
\fancyhead[R]{\small\bfseries\color{violet!85!black}Lớp Toán Thầy Lĩnh -- loptoanthaylinh.vn}
\fancyfoot[C]{\small\color{gray}-- Sổ tay kiến thức: Bài 45. Luyện tập chung phép nhân, phép chia --}

\setlength{\parindent}{0pt}
\setlength{\parskip}{0pt}

\begin{document}

% --- BANNER BÀI HỌC ---
\begin{tcolorbox}[
  enhanced,
  colback=violet!8!white,
  colframe=violet!70!black,
  arc=3mm, boxrule=1.2pt,
  left=3mm, right=3mm, top=2mm, bottom=2mm,
  halign=center
]
  {\LARGE\bfseries\color{violet!90!black}\faBookmark\ BÀI 45: LUYỆN TẬP CHUNG PHÉP NHÂN, PHÉP CHIA}\\[1mm]
  {\itshape\footnotesize\color{black!75} Đại tổng kết 4 bảng cửu chương: Bảng nhân 2, 5 và Bảng chia 2, 5 trên 1 trang A4}
\end{tcolorbox}

\vspace{-2mm}

\noindent
\begin{minipage}[t]{0.49\textwidth}\vspace{0pt}
\begin{tcolorbox}[
  enhanced,
  colback=violet!4!white,
  colframe=violet!75!black,
  arc=2mm, boxrule=0.8pt,
  left=2mm, right=2mm, top=1.8mm, bottom=1.8mm,
  title={\footnotesize\bfseries\color{white}\faKey\ 1. Bốn bảng cửu chương nền tảng}
]
{\footnotesize
\begin{itemize}\setlength{\itemsep}{1pt}\setlength{\leftskip}{-2.5mm}
  \item \textbf{Bảng nhân 2:} Đếm thêm 2 ($2, 4, 6, \dots, 20$).
  \item \textbf{Bảng nhân 5:} Đếm thêm 5 ($5, 10, 15, \dots, 50$).
  \item \textbf{Bảng chia 2:} Chia đôi, nhẩm $2 \times ? = \text{Số bị chia}$.
  \item \textbf{Bảng chia 5:} Chia 5, nhẩm $5 \times ? = \text{Số bị chia}$.
  \item \textbf{Số 1 và số 0 đặc biệt:}\\
  - Nhân với $1$: $a \times 1 = 1 \times a = a$.\\
  - Nhân với $0$: $a \times 0 = 0 \times a = 0$.\\
  - Chia cho $1$: $a : 1 = a$. Số $0 : a = 0$.
\end{itemize}
}
\end{tcolorbox}
\end{minipage}
\hfill
\begin{minipage}[t]{0.49\textwidth}\vspace{0pt}
\begin{tcolorbox}[
  enhanced,
  colback=teal!4!white,
  colframe=teal!75!black,
  arc=2mm, boxrule=0.8pt,
  left=2mm, right=2mm, top=1.8mm, bottom=1.8mm,
  title={\footnotesize\bfseries\color{white}\faSliders\ 2. Kĩ năng tính dãy biểu thức}
]
{\footnotesize
\begin{itemize}\setlength{\itemsep}{1pt}\setlength{\leftskip}{-2.5mm}
  \item \textbf{Biểu thức có phép nhân/chia và cộng/trừ:}\\
  Ta thực hiện \textbf{Nhân, chia trước} $\implies$ \textbf{Cộng, trừ sau}.
  $$\text{Ví dụ: } 2 \times 8 + 15 = 16 + 15 = \mathbf{31}$$
  $$\text{Ví dụ: } 45 - 20 : 5 = 45 - 4 = \mathbf{41}$$
  \item \textbf{Biểu thức chỉ có nhân, chia:}\\
  Thực hiện lần lượt từ \textbf{trái sang phải}:\\
  $2 \times 6 : 4$ (lớp 3) hoặc $20 : 2 \times 0 = 0$.
\end{itemize}
}
\end{tcolorbox}
\end{minipage}

\vspace{-1.5mm}
\begin{tcolorbox}[
  enhanced,
  colback=orange!4!white,
  colframe=orange!80!black,
  arc=2mm, boxrule=0.8pt,
  left=2mm, right=2mm, top=1.5mm, bottom=1.5mm,
  title={\footnotesize\bfseries\color{white}\faSitemap\ 3. Sơ đồ tư duy 4 nhánh: Bộ tứ Bảng cửu chương lớp 2}
]
\centering
\begin{tikzpicture}[
  node distance=0.8cm,
  box/.style={rectangle, rounded corners=2pt, draw=violet!80!black, fill=violet!15, text width=2.4cm, align=center, font=\scriptsize\bfseries, inner sep=3pt},
  branch/.style={rectangle, rounded corners=2pt, draw=#1!80!black, fill=#1!10, text width=2.4cm, align=center, font=\scriptsize, inner sep=2.5pt},
  arr/.style={-{Stealth[scale=0.8]}, thick, #1!80!black}
]
  \node[box] (root) {BẢNG NHÂN\\BẢNG CHIA};
  \node[branch=teal, above left=0.3cm and 0.4cm of root] (b1) {\textbf{BẢNG NHÂN 2}\\Dãy chẵn\\$2 \times 1 \to 2 \times 10$};
  \node[branch=orange, above right=0.3cm and 0.4cm of root] (b2) {\textbf{BẢNG NHÂN 5}\\Tận cùng 0 hoặc 5\\$5 \times 1 \to 5 \times 10$};
  \node[branch=green, below left=0.3cm and 0.4cm of root] (b3) {\textbf{BẢNG CHIA 2}\\Chia đôi\\Ngược nhân 2};
  \node[branch=purple, below right=0.3cm and 0.4cm of root] (b4) {\textbf{BẢNG CHIA 5}\\Một phần năm\\Ngược nhân 5};

  \draw[arr=teal] (root) -- (b1);
  \draw[arr=orange] (root) -- (b2);
  \draw[arr=green] (root) -- (b3);
  \draw[arr=purple] (root) -- (b4);
\end{tikzpicture}
\end{tcolorbox}

\vspace{-1.5mm}
\noindent
\begin{minipage}[t]{0.49\textwidth}\vspace{0pt}
\begin{tcolorbox}[
  enhanced,
  colback=green!4!white,
  colframe=green!60!black,
  arc=2mm, boxrule=0.8pt,
  left=2mm, right=2mm, top=1.8mm, bottom=1.8mm,
  title={\footnotesize\bfseries\color{white}\faCheckCircle\ 4. Ví dụ thực hành chuẩn mực}
]
{\footnotesize
\textbf{1. Tính giá trị biểu thức:}\\
- $5 \times 7 - 12 = 35 - 12 = \mathbf{23}$.\\
- $18 : 2 + 25 = 9 + 25 = \mathbf{34}$.
\textbf{2. So sánh ($>, <, =$):}\\
So sánh: $2 \times 9 \dots 5 \times 4$.\\
Ta có: $2 \times 9 = 18$; $5 \times 4 = 20$.\\
Vì $18 < 20$ nên: $2 \times 9 < 5 \times 4$.
}
\end{tcolorbox}
\end{minipage}
\hfill
\begin{minipage}[t]{0.49\textwidth}\vspace{0pt}
\begin{tcolorbox}[
  enhanced,
  colback=red!4!white,
  colframe=red!75!black,
  arc=2mm, boxrule=0.8pt,
  left=2mm, right=2mm, top=1.8mm, bottom=1.8mm,
  title={\footnotesize\bfseries\color{white}\faExclamationTriangle\ 5. Mẹo tránh bẫy trắc nghiệm}
]
{\footnotesize
\begin{itemize}\setlength{\itemsep}{1pt}\setlength{\leftskip}{-2.5mm}
  \item \textbf{Bẫy thứ tự tính toán:}\\
  $30 - 10 : 2$: Nhiều bạn lấy $30 - 10 = 20$ rồi chia $2 = 10 \implies$ \textbf{sai nặng}!\\
  Phải làm chia trước: $10 : 2 = 5 \implies 30 - 5 = 25$.
  \item Quy tắc vàng: \textbf{"Nhân chia trước, cộng trừ sau"}.
\end{itemize}
}
\end{tcolorbox}
\end{minipage}


\end{document}
